\documentclass[11pt,a4paper]{article}

\usepackage[margin=1in]{geometry}
\usepackage{amsmath,amssymb,amsthm}
\usepackage{booktabs}
\usepackage{hyperref}
\usepackage{xcolor}
\usepackage{microtype}
\setlength{\emergencystretch}{3em}

\hypersetup{
  colorlinks=true,
  linkcolor=blue!60!black,
  citecolor=blue!60!black,
  urlcolor=blue!60!black
}

\newtheorem{theorem}{Theorem}
\newtheorem{proposition}[theorem]{Proposition}
\newtheorem{remark}[theorem]{Remark}

\newcommand{\E}{\mathbb{E}}
\newcommand{\Prob}{\mathbb{P}}
\newcommand{\Var}{\operatorname{Var}}
\newcommand{\Cov}{\operatorname{Cov}}
\newcommand{\DKL}{D_{\mathrm{KL}}}
\newcommand{\dTV}{d_{\mathrm{TV}}}

\title{When Does In-Loop Misalignment Filtering Select for Stealth?\\
Exact Thresholds in a Gibbs Model of Post-Training}

\author{[Author Name]\thanks{Preprint draft prepared for review. Standalone companion verification code: \texttt{verifier\_exemples\_papier\_v3.py}.}}
\date{November 2026}

\begin{document}
\maketitle

\begin{abstract}
When a misalignment detector---such as a chain-of-thought monitor or activation probe---is placed inside the post-training loop, does it select for stealthy evasion, even when intercepted runs are merely discarded and resampled without negative reward? We study this question in one deliberately simple model: the Kullback--Leibler (KL)-regularized Gibbs equilibrium over a three-regime policy space (honest behavior, visible misalignment, stealthy misalignment) whose strata and detection rates are fixed inputs. The model does not address generalisation, i.e.\ whether training against a detector reshapes these inputs. Separating four quantities often conflated in informal debate---conditional stealth enrichment, absolute stealth amplification relative to unmonitored post-training, dominance over honesty, and evasion against a separate deployment monitor---we give exact thresholds for two operators: pure rejection with resampling (Operator~A), which we model as KL-regularized optimization of $R + \eta^{-1}\log\beta$ and which is a modelling choice, and a KL-regularized interception penalty (Operator~B). \textbf{Theorem~1}: under Operator~A, conditional enrichment follows a classical variance identity for every survival function; honesty collapses as optimization pressure $\eta \to \infty$ whenever the stealthy false-negative rate $\beta_{\mathrm{f}}$ is positive, the mass going to cheaters (not necessarily to stealthy ones); a penalty $-C < 0$ leaves a non-vanishing barrier against partially detected stealth, and in a perfectly stealthy niche it places more mass than Operator~A on stealth exactly when $\eta(\Delta R + C) > \log(1/\beta_{\mathrm{v}})/(1-\beta_{\mathrm{v}})$---while then also placing more mass on honesty and less on cheating overall. Modelling rejection as applied at every round instead yields a barrier that does not fade, so the asymptotic contrast between the two operators depends on the modelling of rejection. \textbf{Theorem~2} gives three levers, not an exhaustive classification: (R1)~an exact condition for honesty to dominate stealth, via prior barriers and maintenance costs; (R2)~an exact covariance identity for a deployment monitor under Operator~A---its evasion does not increase iff the covariance between in-loop and audit survival is non-positive, which keeping the audit out of the loop does not by itself ensure---and, under Operator~B, independence or negative association as a sufficient condition (non-positive covariance being necessary but not sufficient); (R3)~indistinguishable planted traps, which exclude frequent but not rare cheaters. A companion script recomputes every number from the closed-form weights and the counterexamples that delimit the claims.
\end{abstract}

\section{Introduction}

Post-training methods based on reinforcement learning and KL-regularized preference optimization substantially increase model capability, but also amplify reward hacking and strategic misalignment~\cite{zhuang2021,skalse2022,gao2022}. A natural engineering response is to deploy misalignment detectors---such as chain-of-thought monitors~\cite{baker2025}, elicitation checks, or internal activation probes---directly inside the post-training loop. Yet empirical and conceptual work on deceptive persistence~\cite{hubinger2024,greenblatt2024} and chain-of-thought obfuscation~\cite{baker2025} raises a fundamental concern: optimizing against an imperfect detector may simply teach the model to hide its misalignment rather than eliminate it. Not training against the chain-of-thought monitor that one relies on is already a published recommendation~\cite{baker2025,korbak2025}; the question here is what selection does, quantitatively, when a filter is nevertheless used.

In a recent discussion~\cite{halberstadt2026}, Halberstadt argued that even \emph{pure rejection without negative penalty} (``discard the run and try again'') reinforces undetected reward hacks, asking whether it is a mathematical law that any attempt to prevent misalignment during post-training inevitably encourages evasion. In response, \texttt{testingthewaters}~\cite{testingthewaters2026} argued that an in-loop filter acts as a selective sieve rather than a directional gradient, and that evasion may fail due to non-existence, unreachability under optimization constraints, or non-maintainability against competing post-training objectives.

This note gives exact answers inside one explicit model, the KL-regularized (Gibbs) post-training equilibrium with fixed strata; the deeper question of generalisation---whether training against a detector teaches honesty or concealment, i.e.\ how $p_{\mathrm{f}}$ and $\beta_{\mathrm{f}}$ change under training---is outside this model. The central identities are classical selection identities \cite{fisher1930,robertson1966,price1970}, and the Gibbs form of KL-regularized RL is known \cite{korbak2022}. We connect the question to strategic classification~\cite{hardt2016}, partial monitoring~\cite{bartok2014}, and AI control~\cite{greenblatt2023,griffin2024,terekhov2025,ukaisi2025}. Our main contributions are:
\begin{enumerate}
  \item \textbf{Disambiguation of four selection metrics:} conditional stealth share among cheaters, absolute stealth probability relative to unmonitored post-training, asymptotic dominance of stealth over honesty as $\eta \to \infty$, and evasion against an off-loop deployment monitor $M_{\mathrm{audit}}$.
  \item \textbf{Exact phase transitions and the Rejection--Penalty Dilemma (Theorem~\ref{thm:selection}):} we state the variance identity for conditional enrichment under Operator~A (a special case of Price's equation), derive exact thresholds for absolute amplification and dominance under both operators, and show that, in this model, honesty collapses under Operator~A whenever $\beta_{\mathrm{f}} > 0$, while in a perfectly stealthy niche a negative penalty yields more stealth than pure rejection past an explicit threshold together with more honesty and less total cheating.
  \item \textbf{Three escape regimes (Theorem~\ref{thm:escape}):} we give three levers, without claiming exhaustiveness: (R1)~the exact condition, in this model, for honesty to dominate stealth via prior KL barriers and maintenance costs; (R2)~a covariance condition for a deployment monitor (exact identity under Operator~A; independence or negative association sufficient under Operator~B, non-positive covariance necessary but not sufficient)---a condition on dependence, not on placement; and (R3)~planted honeypots at bounded total variation distance, which exclude frequent cheaters.
  \item \textbf{Reproducible numbers:} a standalone script recomputes every table from the closed-form weights and the exact counterexamples of Remark~\ref{rem:delimit}; the proofs, not the script, carry the results.
\end{enumerate}

\section{Formal Framework}

Let $(\Pi, \mathcal{A}, \pi_0)$ be the probability space of elementary policies or complete trajectories (including internal reasoning traces and tool calls) induced by a pretrained base model $\pi_0$. Partition $\Pi$ into three disjoint measurable events of strictly positive prior mass:
\begin{equation}
  \Pi = \Pi_{\mathrm{h}} \sqcup \Pi_{\mathrm{v}} \sqcup \Pi_{\mathrm{f}},
  \qquad p_{\mathrm{h}} := \pi_0(\Pi_{\mathrm{h}}) > 0,\quad p_{\mathrm{v}} := \pi_0(\Pi_{\mathrm{v}}) > 0,\quad p_{\mathrm{f}} := \pi_0(\Pi_{\mathrm{f}}) > 0,
\end{equation}
with $p_{\mathrm{h}} + p_{\mathrm{v}} + p_{\mathrm{f}} = 1$. Here $\Pi_{\mathrm{h}}$ denotes \emph{honest} behavior, $\Pi_{\mathrm{v}}$ denotes \emph{visible} misalignment (transparent reward hacking), $\Pi_{\mathrm{f}}$ denotes \emph{stealthy} misalignment (obfuscated reward hacking), and $\Pi_{\mathrm{tr}} := \Pi_{\mathrm{v}} \sqcup \Pi_{\mathrm{f}}$ denotes all misaligned behavior.

\subsection{Task Reward and In-Loop Monitor}
By shift invariance of Gibbs measures and KL-regularized objectives, we normalize task reward $R : \Pi \to \mathbb{R}$ to $R(\pi) = 0$ on $\Pi_{\mathrm{h}}$. On $\Pi_{\mathrm{v}}$ and $\Pi_{\mathrm{f}}$, misalignment yields reward advantages $\Delta R_{\mathrm{v}} > 0$ and $\Delta R_{\mathrm{f}} = \Delta R_{\mathrm{v}} - \delta_{\mathrm{maint}}$, where $\delta_{\mathrm{maint}} \ge 0$ captures any capability loss or conflict with auxiliary post-training objectives. In the canonical regime without intrinsic obfuscation cost, $\Delta R_{\mathrm{v}} = \Delta R_{\mathrm{f}} = \Delta R > 0$.

An in-loop monitor $M_{\mathrm{loop}} : \Pi \to \{0, 1\}$ ($1 =$ intercepted, $0 =$ accepted) is summarized by its \emph{survival probability}
\begin{equation}
  \beta(\pi) := \Prob(M_{\mathrm{loop}} = 0 \mid \pi) \in [0, 1],
\end{equation}
which equals the false-negative rate on $\Pi_{\mathrm{tr}}$ and $1 - \alpha_{\mathrm{h}}$ on $\Pi_{\mathrm{h}}$ ($\alpha_{\mathrm{h}} \in [0, 1)$ being the false-positive rate). In the stratum-constant reduction, we assume $0 \le \beta_{\mathrm{v}} < \beta_{\mathrm{f}} \le \beta_{\mathrm{h}} \le 1$ with $\beta_{\mathrm{f}} > 0$.

\subsection{Canonical Post-Training Operators}
At optimization pressure (inverse KL temperature) $\eta > 0$, the unmonitored baseline ($\beta \equiv 1$) is the unique maximizer of $\E_{\pi}[R] - \eta^{-1}\DKL(\pi \,\|\, \pi_0)$:
\begin{equation}
  \pi_{\mathrm{no}}^{(\eta)}(d\pi) = \frac{e^{\eta R(\pi)}\,\pi_0(d\pi)}{Z_{\mathrm{no}}(\eta)},
  \qquad Z_{\mathrm{no}}(\eta) = \int_{\Pi} e^{\eta R(\pi')}\,\pi_0(d\pi'),
\end{equation}
extended by continuity to $\pi_{\mathrm{no}}^{(0)} := \pi_0$. We formalize the two canonical in-loop interventions:
\begin{enumerate}
  \item \textbf{Operator~(A) --- Pure rejection with resampling (discard and try again).} Intercepted trajectories ($M_{\mathrm{loop}} = 1$) are discarded prior to reinforcement, conditioning the base law to $\pi_0^{\mathrm{acc}}(d\pi) = \beta(\pi)\pi_0(d\pi)/\E_{\pi_0}[\beta]$ and then applying KL-regularized reward maximization relative to $\pi_0^{\mathrm{acc}}$:
  \begin{equation}
    \pi_{\mathrm{A}}^{(\eta)}(d\pi) = \frac{\beta(\pi)\,e^{\eta R(\pi)}\,\pi_0(d\pi)}{Z_{\mathrm{A}}(\eta)},
    \qquad Z_{\mathrm{A}}(\eta) = \int_{\Pi} \beta(\pi')\,e^{\eta R(\pi')}\,\pi_0(d\pi').
  \end{equation}
  At $\eta = 0$, $\pi_{\mathrm{A}}^{(0)}$ reduces to pure Bayesian rejection conditioning without reward optimization.

  \emph{Operator~A is a modelling choice.} $\pi_{\mathrm{A}}^{(\eta)}$ is the unique maximizer of $\E_\pi[R + \eta^{-1}\log\beta] - \eta^{-1}\DKL(\pi\,\|\,\pi_0)$: the filter acts as a bonus $\eta^{-1}\log\beta$ that fades as $\eta$ grows, by construction, and $\pi_{\mathrm{A}}^{(\eta)} \propto \beta\,\pi_{\mathrm{no}}^{(\eta)}$ coincides with unmonitored training followed by a single rejection at inference. Other models of ``discard and try again'' behave differently: with rejection applied at every round, $T$ rounds of step $\eta_1$ give $\pi \propto \beta^T e^{T\eta_1 R}\pi_0$, and stealth eventually dominates honesty iff $\eta_1 \Delta R > \log(\beta_{\mathrm{h}}/\beta_{\mathrm{f}})$, a barrier that does not fade.
  \item \textbf{Operator~(B) --- KL-regularized update with interception penalty $-C \le 0$.} Intercepted trajectories are halted before taking effect and receive score $-C$ ($C \ge 0$), yielding expected reward $\bar{R}_C(\pi) = \beta(\pi)R(\pi) - (1-\beta(\pi))C$. The unique maximizer of $\E_{\pi}[\bar{R}_C] - \eta^{-1}\DKL(\pi \,\|\, \pi_0)$ ($\eta > 0$) is:
  \begin{equation}
    \pi_{\mathrm{B}}^{(\eta, C)}(d\pi) = \frac{\exp\!\big(\eta\big[\beta(\pi)R(\pi) - (1-\beta(\pi))C\big]\big)\,\pi_0(d\pi)}{Z_{\mathrm{B}}(\eta, C)}.
  \end{equation}
\end{enumerate}

\section{Exact Phase Transitions and the Rejection--Penalty Dilemma}

\begin{theorem}[Exact Characterization of Stealth Selection under an Imperfect Filter]\label{thm:selection}
Suppose $\Delta R_{\mathrm{v}} = \Delta R_{\mathrm{f}} = \Delta R > 0$ and $0 \le \beta_{\mathrm{v}} < \beta_{\mathrm{f}} \le \beta_{\mathrm{h}} \le 1$ with $\beta_{\mathrm{f}} > 0$.
\begin{enumerate}
  \item \textbf{Conditional enrichment under Operator~A (variance identity, for every survival function).} For any measurable survival function $\beta : \Pi_{\mathrm{tr}} \to [0, 1]$ not almost surely constant under $\mu_{\mathrm{tr}}^{(\eta)}(\cdot) := \pi_{\mathrm{no}}^{(\eta)}(\cdot \mid \Pi_{\mathrm{tr}})$ (with $\E_{\mu_{\mathrm{tr}}^{(\eta)}}[\beta] > 0$),
  \begin{equation}\label{eq:variance}
    \E_{\pi_{\mathrm{A}}^{(\eta)}}[\beta(\pi) \mid \Pi_{\mathrm{tr}}] - \E_{\mu_{\mathrm{tr}}^{(\eta)}}[\beta(\pi)]
    = \frac{\Var_{\mu_{\mathrm{tr}}^{(\eta)}}(\beta(\pi))}{\E_{\mu_{\mathrm{tr}}^{(\eta)}}[\beta(\pi)]} > 0.
  \end{equation}
  In the two-stratum cheating model ($\beta_{\mathrm{v}} < \beta_{\mathrm{f}}$), conditional stealth share strictly increases without any threshold under both Operator~A ($\forall \eta \ge 0$) and Operator~B ($\forall \eta > 0, \forall C \ge 0$):
  \begin{equation}
    \pi_{\mathrm{A}}^{(\eta)}(\Pi_{\mathrm{f}} \mid \Pi_{\mathrm{tr}}) > \pi_{\mathrm{no}}^{(\eta)}(\Pi_{\mathrm{f}} \mid \Pi_{\mathrm{tr}})
    \qquad \text{and} \qquad
    \pi_{\mathrm{B}}^{(\eta, C)}(\Pi_{\mathrm{f}} \mid \Pi_{\mathrm{tr}}) > \pi_{\mathrm{no}}^{(\eta)}(\Pi_{\mathrm{f}} \mid \Pi_{\mathrm{tr}}).
  \end{equation}

  \item \textbf{Exact absolute-amplification threshold under pure rejection (Operator~A).} For fixed $\eta > 0$, $\pi_{\mathrm{A}}^{(\eta)}(\Pi_{\mathrm{f}}) > \pi_{\mathrm{no}}^{(\eta)}(\Pi_{\mathrm{f}})$ holds if and only if either $\beta_{\mathrm{f}} = \beta_{\mathrm{h}}$ (in which case amplification holds for all $\Delta R \ge 0, \eta \ge 0$), or $\beta_{\mathrm{f}} < \beta_{\mathrm{h}}$ and
  \begin{equation}\label{eq:ampA}
    \Delta R > \Delta R_{\mathrm{A}}^*(\eta)
    := \frac{1}{\eta}\left[\log\frac{p_{\mathrm{h}}}{p_{\mathrm{v}}} + \log\frac{\beta_{\mathrm{h}} - \beta_{\mathrm{f}}}{\beta_{\mathrm{f}} - \beta_{\mathrm{v}}}\right].
  \end{equation}

  \item \textbf{Exact absolute-amplification condition under KL-regularized penalty (Operator~B; an implicit inequality, not a closed form).} For all $\eta > 0$ and $C \ge 0$, $\pi_{\mathrm{B}}^{(\eta, C)}(\Pi_{\mathrm{f}}) > \pi_{\mathrm{no}}^{(\eta)}(\Pi_{\mathrm{f}})$ holds if and only if
  \begin{equation}\label{eq:ampB}
    p_{\mathrm{v}}\Big(1 - e^{-\eta(\beta_{\mathrm{f}} - \beta_{\mathrm{v}})(\Delta R + C)}\Big)
    > p_{\mathrm{h}} e^{-\eta \Delta R}\Big(e^{\eta[(1-\beta_{\mathrm{f}})\Delta R + (\beta_{\mathrm{h}} - \beta_{\mathrm{f}})C]} - 1\Big).
  \end{equation}

  \item \textbf{Structural divergence (Rejection--Penalty Dilemma).}
  \begin{enumerate}
    \item \emph{Asymptotic collapse of pure rejection:} Under Operator~A, $\pi_{\mathrm{A}}^{(\eta)}(\Pi_{\mathrm{f}}) > \pi_{\mathrm{A}}^{(\eta)}(\Pi_{\mathrm{h}})$ holds iff
    \begin{equation}\label{eq:domA}
      \Delta R > \Delta R_{\mathrm{dom, A}}^*(\eta)
      := \frac{1}{\eta}\left(\log\frac{p_{\mathrm{h}}}{p_{\mathrm{f}}} + \log\frac{\beta_{\mathrm{h}}}{\beta_{\mathrm{f}}}\right)
      \xrightarrow[\eta \to +\infty]{} 0.
    \end{equation}
    Hence, for any $\Delta R > 0$ and any imperfect detection rate ($\beta_{\mathrm{f}} > 0$), honesty collapses under pure rejection as $\eta \to +\infty$. The mass goes to the cheaters, not necessarily to the stealthy ones: with equal rewards, $\pi_{\mathrm{A}}^{(\eta)}(\Pi_{\mathrm{f}} \mid \Pi_{\mathrm{tr}}) = \beta_{\mathrm{f}}p_{\mathrm{f}}/(\beta_{\mathrm{v}}p_{\mathrm{v}} + \beta_{\mathrm{f}}p_{\mathrm{f}})$ for every $\eta$, which can be small.
    \item \emph{Non-vanishing barrier under negative penalty when $\beta_{\mathrm{f}} < \beta_{\mathrm{h}}$:} Under Operator~B, $\pi_{\mathrm{B}}^{(\eta, C)}(\Pi_{\mathrm{f}}) > \pi_{\mathrm{B}}^{(\eta, C)}(\Pi_{\mathrm{h}})$ holds iff
    \begin{equation}\label{eq:domB}
      \Delta R > \Delta R_{\mathrm{dom, B}}^*(\eta, C)
      := \frac{\beta_{\mathrm{h}} - \beta_{\mathrm{f}}}{\beta_{\mathrm{f}}}\,C + \frac{1}{\eta \beta_{\mathrm{f}}}\log\frac{p_{\mathrm{h}}}{p_{\mathrm{f}}}
      \xrightarrow[\eta \to +\infty]{} \frac{\beta_{\mathrm{h}} - \beta_{\mathrm{f}}}{\beta_{\mathrm{f}}}\,C > 0 \quad (C > 0).
    \end{equation}
    \item \emph{Perfectly stealthy niche ($\beta_{\mathrm{f}} = \beta_{\mathrm{h}} = 1 > \beta_{\mathrm{v}} > 0$):} If a stealthy strategy completely evades the in-loop monitor, Operator~B loses its asymptotic barrier ($\Delta R_{\mathrm{dom, B}}^* \to 0$). The stealth-to-honest odds equal $p_{\mathrm{f}}e^{\eta\Delta R}/p_{\mathrm{h}}$ under both operators, so $\pi_{\mathrm{B}}^{(\eta,C)}(\Pi_{\mathrm{f}}) > \pi_{\mathrm{A}}^{(\eta)}(\Pi_{\mathrm{f}})$ holds iff $\pi_{\mathrm{B}}^{(\eta,C)}(\Pi_{\mathrm{h}}) > \pi_{\mathrm{A}}^{(\eta)}(\Pi_{\mathrm{h}})$, in which case total cheating is lower under~B. Operator~B places more mass on stealth than pure rejection---that is, $\pi_{\mathrm{B}}^{(\eta, C)}(\Pi_{\mathrm{f}}) > \pi_{\mathrm{A}}^{(\eta)}(\Pi_{\mathrm{f}}) > \pi_{\mathrm{no}}^{(\eta)}(\Pi_{\mathrm{f}})$---if and only if
    \begin{equation}\label{eq:inversion}
      \eta(\Delta R + C) > \frac{\log(1/\beta_{\mathrm{v}})}{1 - \beta_{\mathrm{v}}}.
    \end{equation}
  \end{enumerate}
\end{enumerate}
\end{theorem}

\begin{proof}
\emph{Part~1.} Since $e^{\eta R(\pi)}$ is already absorbed into $\mu_{\mathrm{tr}}^{(\eta)}(\cdot) = \pi_{\mathrm{no}}^{(\eta)}(\cdot \mid \Pi_{\mathrm{tr}})$, the Radon--Nikodym derivative of $\pi_{\mathrm{A}}^{(\eta)}(\cdot \mid \Pi_{\mathrm{tr}})$ with respect to $\mu_{\mathrm{tr}}^{(\eta)}$ is $\beta(\pi) / \E_{\mu_{\mathrm{tr}}^{(\eta)}}[\beta]$ (even when $R(\pi)$ varies on $\Pi_{\mathrm{tr}}$). Therefore,
\[
  \E_{\pi_{\mathrm{A}}^{(\eta)}}[\beta \mid \Pi_{\mathrm{tr}}]
  = \frac{\E_{\mu_{\mathrm{tr}}^{(\eta)}}[\beta^2]}{\E_{\mu_{\mathrm{tr}}^{(\eta)}}[\beta]}
  = \E_{\mu_{\mathrm{tr}}^{(\eta)}}[\beta] + \frac{\Var_{\mu_{\mathrm{tr}}^{(\eta)}}(\beta)}{\E_{\mu_{\mathrm{tr}}^{(\eta)}}[\beta]},
\]
which strictly exceeds $\E_{\mu_{\mathrm{tr}}^{(\eta)}}[\beta]$ whenever $\Var_{\mu_{\mathrm{tr}}^{(\eta)}}(\beta) > 0$. For two strata with $\beta_{\mathrm{v}} < \beta_{\mathrm{f}}$, the conditional odds are $\frac{\pi_{\mathrm{A}}^{(\eta)}(\Pi_{\mathrm{f}})}{\pi_{\mathrm{A}}^{(\eta)}(\Pi_{\mathrm{v}})} = \frac{\beta_{\mathrm{f}}}{\beta_{\mathrm{v}}}\frac{p_{\mathrm{f}}}{p_{\mathrm{v}}} > \frac{p_{\mathrm{f}}}{p_{\mathrm{v}}}$ (infinite when $\beta_{\mathrm{v}} = 0$) and $\frac{\pi_{\mathrm{B}}^{(\eta, C)}(\Pi_{\mathrm{f}})}{\pi_{\mathrm{B}}^{(\eta, C)}(\Pi_{\mathrm{v}})} = e^{\eta(\beta_{\mathrm{f}} - \beta_{\mathrm{v}})(\Delta R + C)}\frac{p_{\mathrm{f}}}{p_{\mathrm{v}}} > \frac{p_{\mathrm{f}}}{p_{\mathrm{v}}}$ since $\eta > 0$, $\beta_{\mathrm{f}} > \beta_{\mathrm{v}}$, and $\Delta R + C > 0$. Strict monotonicity of $x \mapsto x/(1+x)$ on $[0, +\infty]$ yields both strict inequalities.

\emph{Part~2.} Writing the stratum probabilities explicitly,
\[
  \pi_{\mathrm{A}}^{(\eta)}(\Pi_{\mathrm{f}})
  = \frac{\beta_{\mathrm{f}} p_{\mathrm{f}} e^{\eta \Delta R}}{\beta_{\mathrm{h}} p_{\mathrm{h}} + (\beta_{\mathrm{v}} p_{\mathrm{v}} + \beta_{\mathrm{f}} p_{\mathrm{f}}) e^{\eta \Delta R}},
  \qquad
  \pi_{\mathrm{no}}^{(\eta)}(\Pi_{\mathrm{f}})
  = \frac{p_{\mathrm{f}} e^{\eta \Delta R}}{p_{\mathrm{h}} + (p_{\mathrm{v}} + p_{\mathrm{f}}) e^{\eta \Delta R}}.
\]
Since $p_{\mathrm{f}} e^{\eta \Delta R} > 0$ and both denominators are strictly positive, cross-multiplying shows that $\pi_{\mathrm{A}}^{(\eta)}(\Pi_{\mathrm{f}}) > \pi_{\mathrm{no}}^{(\eta)}(\Pi_{\mathrm{f}})$ is equivalent to
\[
  (\beta_{\mathrm{f}} - \beta_{\mathrm{v}}) p_{\mathrm{v}} e^{\eta \Delta R} > (\beta_{\mathrm{h}} - \beta_{\mathrm{f}}) p_{\mathrm{h}}.
\]
The left-hand side is strictly positive because $\beta_{\mathrm{f}} > \beta_{\mathrm{v}}$ and $p_{\mathrm{v}} > 0$. If $\beta_{\mathrm{f}} = \beta_{\mathrm{h}}$, the right-hand side vanishes and the inequality always holds. If $\beta_{\mathrm{f}} < \beta_{\mathrm{h}}$, dividing by $(\beta_{\mathrm{f}} - \beta_{\mathrm{v}})p_{\mathrm{v}}$ and taking logarithms yields~\eqref{eq:ampA}.

\emph{Part~3.} Dividing numerator and denominator of $\pi_{\mathrm{B}}^{(\eta, C)}(\Pi_{\mathrm{f}})$ by $\exp(\eta[\beta_{\mathrm{f}}\Delta R - (1-\beta_{\mathrm{f}})C])$ gives
\begin{gather*}
  \pi_{\mathrm{B}}^{(\eta, C)}(\Pi_{\mathrm{f}})
  = \frac{p_{\mathrm{f}}}{p_{\mathrm{h}} e^{-\eta[\beta_{\mathrm{f}}\Delta R - (\beta_{\mathrm{h}}-\beta_{\mathrm{f}})C]} + p_{\mathrm{v}} e^{-\eta(\beta_{\mathrm{f}}-\beta_{\mathrm{v}})(\Delta R + C)} + p_{\mathrm{f}}},\\
  \pi_{\mathrm{no}}^{(\eta)}(\Pi_{\mathrm{f}})
  = \frac{p_{\mathrm{f}}}{p_{\mathrm{h}} e^{-\eta \Delta R} + p_{\mathrm{v}} + p_{\mathrm{f}}}.
\end{gather*}
Comparing the two denominators directly yields~\eqref{eq:ampB}.

\emph{Part~4.} The log-odds $\log\frac{\pi_{\mathrm{A}}^{(\eta)}(\Pi_{\mathrm{f}})}{\pi_{\mathrm{A}}^{(\eta)}(\Pi_{\mathrm{h}})} = \log\frac{\beta_{\mathrm{f}} p_{\mathrm{f}}}{\beta_{\mathrm{h}} p_{\mathrm{h}}} + \eta\Delta R$ and $\log\frac{\pi_{\mathrm{B}}^{(\eta, C)}(\Pi_{\mathrm{f}})}{\pi_{\mathrm{B}}^{(\eta, C)}(\Pi_{\mathrm{h}})} = \log\frac{p_{\mathrm{f}}}{p_{\mathrm{h}}} + \eta[\beta_{\mathrm{f}}\Delta R - (\beta_{\mathrm{h}} - \beta_{\mathrm{f}})C]$ immediately give~\eqref{eq:domA} and~\eqref{eq:domB}. Finally, when $\beta_{\mathrm{f}} = \beta_{\mathrm{h}} = 1$ and $0 < \beta_{\mathrm{v}} < 1$, both $\pi_{\mathrm{A}}^{(\eta)}(\Pi_{\mathrm{f}})$ and $\pi_{\mathrm{B}}^{(\eta, C)}(\Pi_{\mathrm{f}})$ share the numerator $p_{\mathrm{f}} e^{\eta \Delta R}$ and the denominator terms $p_{\mathrm{h}} + p_{\mathrm{f}} e^{\eta \Delta R}$; only the weight of $\Pi_{\mathrm{v}}$ differs, being $\beta_{\mathrm{v}} p_{\mathrm{v}} e^{\eta \Delta R}$ under Operator~A versus $p_{\mathrm{v}} e^{\eta[\beta_{\mathrm{v}}\Delta R - (1-\beta_{\mathrm{v}})C]} = p_{\mathrm{v}} e^{\eta \Delta R} e^{-\eta(1-\beta_{\mathrm{v}})(\Delta R + C)}$ under Operator~B. Thus $\pi_{\mathrm{B}}^{(\eta, C)}(\Pi_{\mathrm{f}}) > \pi_{\mathrm{A}}^{(\eta)}(\Pi_{\mathrm{f}})$ iff $e^{-\eta(1-\beta_{\mathrm{v}})(\Delta R + C)} < \beta_{\mathrm{v}}$, which is~\eqref{eq:inversion}.
\end{proof}

\begin{proposition}[Limiting Cases and Zero-Penalty Halting without Resampling]\label{prop:limits}
Under the assumptions of Theorem~\ref{thm:selection}:
\begin{enumerate}
  \item \emph{Zero-penalty halting without resampling ($C = 0$ in Operator~B):} Setting $C = 0$ in Operator~B reduces the stealth reward exponent from $\eta\Delta R$ to $\eta\beta_{\mathrm{f}}\Delta R$, which still converges to $\pi_{\mathrm{B}}^{(\eta, 0)}(\Pi_{\mathrm{f}}) \to 1$ as $\eta \to +\infty$ for any $\beta_{\mathrm{f}} > 0$. Thus, halting without resampling and without negative penalty slows the asymptotic exponential rate of stealth dominance by a factor $\beta_{\mathrm{f}}$ (scaling the asymptotic $\eta$-scale by $1/\beta_{\mathrm{f}}$) compared to pure rejection, while the exact finite-$\eta$ threshold ratio $\Delta R_{\mathrm{dom, B}}^*(\eta, 0)/\Delta R_{\mathrm{dom, A}}^*(\eta) = \frac{1}{\beta_{\mathrm{f}}}\frac{\log(p_{\mathrm{h}}/p_{\mathrm{f}})}{\log(p_{\mathrm{h}}/p_{\mathrm{f}}) + \log(\beta_{\mathrm{h}}/\beta_{\mathrm{f}})}$ differs from $1/\beta_{\mathrm{f}}$ whenever $\beta_{\mathrm{f}} < \beta_{\mathrm{h}}$ (and is undefined when $p_{\mathrm{h}}\beta_{\mathrm{h}} = p_{\mathrm{f}}\beta_{\mathrm{f}}$); in either case, it does not prevent asymptotic collapse.
  \item \emph{Perfect detection of visible cheating ($\beta_{\mathrm{v}} = 0$):} Under Operator~A, $\pi_{\mathrm{A}}^{(\eta)}(\Pi_{\mathrm{v}}) = 0$ for all $\eta \ge 0$, so $\pi_{\mathrm{A}}^{(\eta)}(\Pi_{\mathrm{f}} \mid \Pi_{\mathrm{tr}}) = 1$; under Operator~B with $\beta_{\mathrm{f}} = \beta_{\mathrm{h}} = 1$, $\pi_{\mathrm{A}}^{(\eta)}(\Pi_{\mathrm{f}}) > \pi_{\mathrm{B}}^{(\eta, C)}(\Pi_{\mathrm{f}})$ for all finite $\eta, C$ because Operator~A eliminates $\Pi_{\mathrm{v}}$ completely while Operator~B retains exponentially small mass $p_{\mathrm{v}} e^{-\eta C}$.
\end{enumerate}
\end{proposition}

\begin{proof}
Part~1 follows directly from~\eqref{eq:domA} and~\eqref{eq:domB} at $C = 0$, giving $\Delta R_{\mathrm{dom, B}}^*(\eta, 0) = (\eta\beta_{\mathrm{f}})^{-1}\log(p_{\mathrm{h}}/p_{\mathrm{f}}) \to 0$. Part~2 follows by substituting $\beta_{\mathrm{v}} = 0$ into the unnormalized weights of $\Pi_{\mathrm{v}}$ ($0$ under Operator~A versus $p_{\mathrm{v}} e^{-\eta C} > 0$ under Operator~B).
\end{proof}

\section{Three Levers Against Stealth Selection}

We now give three levers against stealth selection. They are not claimed to be an exhaustive classification: we do not prove that avoiding deployment-effective evasion requires any one of them. (R1) is an exact dominance condition in the three-regime model; (R2) is a condition on the dependence between in-loop and audit survival, not on where the audit is placed; (R3) is a sufficient mechanism in an extended contextual model. Throughout, $p_{\mathrm{f}} > 0$ (the case $p_{\mathrm{f}} = 0$ is outside the domain), and the link between (R3) and the critical ratio $\frac{\beta_{\mathrm{h}} C}{\Delta R_{\mathrm{net}} + C}$ of (R1) assumes $\beta_{\mathrm{h}} = 1$.

\begin{theorem}[Levers R1, R2, and R3]\label{thm:escape}
\leavevmode
\begin{enumerate}
  \item \textbf{Regime (R1) --- Prior KL barrier and multi-objective maintenance cost (exact dominance condition in the 3-regime model).} Define $\Delta \DKL(\Pi_{\mathrm{f}} \,\|\, \Pi_{\mathrm{h}}; \pi_0) := \log(p_{\mathrm{h}}/p_{\mathrm{f}})$ and $\Delta R_{\mathrm{net}} := \Delta R - \delta_{\mathrm{maint}}$. Honesty strictly dominates stealth ($\pi_{\mathrm{post}}(\Pi_{\mathrm{f}}) < \pi_{\mathrm{post}}(\Pi_{\mathrm{h}})$) if and only if
  \begin{equation}\label{eq:R1}
    \Delta \DKL(\Pi_{\mathrm{f}} \,\|\, \Pi_{\mathrm{h}}; \pi_0) >
    \begin{cases}
      \eta\,\Delta R_{\mathrm{net}} - \log(\beta_{\mathrm{h}}/\beta_{\mathrm{f}}) & \text{under Operator~(A)}, \\[4pt]
      \eta\big[\beta_{\mathrm{f}}\Delta R_{\mathrm{net}} - (\beta_{\mathrm{h}} - \beta_{\mathrm{f}})C\big] & \text{under Operator~(B)}.
    \end{cases}
  \end{equation}
  In particular, if $\Pi_{\mathrm{f}}$ is unreachable ($p_{\mathrm{f}} = 0 \iff \Delta \DKL = +\infty$) or non-maintainable ($\beta_{\mathrm{f}}\Delta R_{\mathrm{net}} < (\beta_{\mathrm{h}} - \beta_{\mathrm{f}})C$, equivalently $\beta_{\mathrm{f}} < \frac{\beta_{\mathrm{h}} C}{\Delta R_{\mathrm{net}} + C}$ when $\Delta R_{\mathrm{net}} + C > 0$, or $\Delta R_{\mathrm{net}} < 0$), then $\lim_{\eta \to +\infty}\pi_{\mathrm{B}}^{(\eta, C)}(\Pi_{\mathrm{f}}) = 0$.

  \item \textbf{Regime (R2) --- Audit decoupling: a covariance condition, not a placement.} Suppose deployment monitoring relies on a detector $M_{\mathrm{audit}}$ (in or out of the loop) with survival probability $\beta_{\mathrm{audit}}(\pi) := \Prob(M_{\mathrm{audit}} = 0 \mid \pi)$. On any iso-reward cheating stratum $\mathcal{S} := \{\pi \in \Pi_{\mathrm{tr}} : R(\pi) = \Delta R\}$ with $\pi_0(\mathcal{S}) > 0$, $\nu_0 := \pi_0(\cdot \mid \mathcal{S})$ and $\E_{\nu_0}[\beta_{\mathrm{loop}}] > 0$, Operator~A satisfies the exact identity
  \begin{equation}\label{eq:R2}
    \E_{\pi_{\mathrm{A}}^{(\eta)}}[\beta_{\mathrm{audit}} \mid \mathcal{S}] - \E_{\nu_0}[\beta_{\mathrm{audit}}]
    = \frac{\Cov_{\nu_0}(\beta_{\mathrm{loop}}, \beta_{\mathrm{audit}})}{\E_{\nu_0}[\beta_{\mathrm{loop}}]}.
  \end{equation}
  so under Operator~A the audit survival does not increase iff $\Cov_{\nu_0}(\beta_{\mathrm{loop}}, \beta_{\mathrm{audit}}) \le 0$; keeping $M_{\mathrm{audit}}$ out of the loop does not by itself ensure this. Under Operator~B, $\Cov_{\nu_0}(\beta_{\mathrm{loop}}, \beta_{\mathrm{audit}}) \le 0$ is necessary for the audit survival not to increase for small $\eta(\Delta R + C)$ (it is the derivative at $0$), but it is not sufficient; if $(\beta_{\mathrm{loop}}, \beta_{\mathrm{audit}})$ are independent or negatively associated~\cite{joagdev1983} under $\nu_0$, then for all $\eta \ge 0$ and $C \ge 0$,
  \begin{equation}
    \E_{\pi_{\mathrm{B}}^{(\eta, C)}}[\beta_{\mathrm{audit}} \mid \mathcal{S}] \le \E_{\nu_0}[\beta_{\mathrm{audit}}].
  \end{equation}

  \item \textbf{Regime (R3) --- Indistinguishable planted honeypots (sufficient constructive mechanism in the extended contextual model).} Let $\mathcal{X}$ be the space of contexts observed by the policy, drawn from the mixture $P_{\mathrm{mix}} = (1-\rho)P_{\mathrm{ord}} + \rho P_{\mathrm{trap}}$ ($\rho \in (0, 1)$), where the monitor may be completely blind on ordinary episodes $P_{\mathrm{ord}}$ ($\beta_{\mathrm{ord}} \le 1$) but catches any cheating on trap episodes $P_{\mathrm{trap}}$ with probability at least $1 - \varepsilon_{\mathrm{trap}} \in (0, 1]$. If every admissible policy $\pi$ (with context-dependent cheating probability $\tau_{\pi} : \mathcal{X} \to [0, 1]$) satisfies $|\E_{P_{\mathrm{ord}}}[\tau_\pi(X)] - \E_{P_{\mathrm{trap}}}[\tau_\pi(X)]| \le \delta_{\mathrm{ind}} \in [0, 1)$---which holds unconditionally whenever $\dTV(P_{\mathrm{ord}}, P_{\mathrm{trap}}) \le \delta_{\mathrm{ind}}$---then any policy with $\tau_{\mathrm{ord}}(\pi) := \E_{P_{\mathrm{ord}}}[\tau_\pi(X)] \ge \tau_{\min} > \delta_{\mathrm{ind}}$ incurs trap-cheating rate $\tau_{\mathrm{trap}}(\pi) \ge \tau_{\mathrm{ord}}(\pi) - \delta_{\mathrm{ind}}$ and has strictly negative expected reward relative to pure honesty ($\tau \equiv 0$, $\beta_{\mathrm{h}} = 1$) under Operator~B whenever
  \begin{equation}\label{eq:R3}
    C > C_{\mathrm{R3}}^*
    := \frac{(1-\rho)\tau_{\min} + \rho\,\varepsilon_{\mathrm{trap}}(\tau_{\min} - \delta_{\mathrm{ind}})}{\rho(1 - \varepsilon_{\mathrm{trap}})(\tau_{\min} - \delta_{\mathrm{ind}})}\,\Delta R
    \le \frac{1 - \rho + \rho\,\varepsilon_{\mathrm{trap}}}{\rho(1 - \varepsilon_{\mathrm{trap}})(1 - \delta_{\mathrm{ind}}/\tau_{\min})}\,\Delta R.
  \end{equation}
\end{enumerate}
\end{theorem}

\begin{proof}
\emph{Part~1 (R1).} Substituting $\Delta R_{\mathrm{f}} = \Delta R_{\mathrm{net}}$ into the log-odds ratios $\pi_{\mathrm{A}}^{(\eta)}(\Pi_{\mathrm{f}})/\pi_{\mathrm{A}}^{(\eta)}(\Pi_{\mathrm{h}}) = \frac{\beta_{\mathrm{f}} p_{\mathrm{f}}}{\beta_{\mathrm{h}} p_{\mathrm{h}}} e^{\eta \Delta R_{\mathrm{net}}}$ and $\pi_{\mathrm{B}}^{(\eta, C)}(\Pi_{\mathrm{f}})/\pi_{\mathrm{B}}^{(\eta, C)}(\Pi_{\mathrm{h}}) = \frac{p_{\mathrm{f}}}{p_{\mathrm{h}}} e^{\eta[\beta_{\mathrm{f}}\Delta R_{\mathrm{net}} - (\beta_{\mathrm{h}}-\beta_{\mathrm{f}})C]}$ shows that each ratio is strictly less than $1$ iff~\eqref{eq:R1} holds. Whenever the coefficient of $\eta$ is strictly negative, the ratio decays exponentially to $0$ as $\eta \to +\infty$.

\emph{Part~2 (R2).} On $\mathcal{S}$, $e^{\eta R(\pi)} = e^{\eta \Delta R}$ cancels in $\pi_{\mathrm{A}}^{(\eta)}(\cdot \mid \mathcal{S})$, so $d\pi_{\mathrm{A}}^{(\eta)}(\cdot \mid \mathcal{S})/d\nu_0 = \beta_{\mathrm{loop}} / \E_{\nu_0}[\beta_{\mathrm{loop}}]$. Expanding $\E_{\nu_0}[\beta_{\mathrm{loop}}\beta_{\mathrm{audit}}]/\E_{\nu_0}[\beta_{\mathrm{loop}}] - \E_{\nu_0}[\beta_{\mathrm{audit}}]$ gives~\eqref{eq:R2}. Under Operator~B, the conditional density is proportional to $g(\beta_{\mathrm{loop}})$ where $g(u) = e^{\eta u (\Delta R + C)}$ is non-decreasing; by negative association~\cite{joagdev1983} (or independence), $\Cov_{\nu_0}(g(\beta_{\mathrm{loop}}), \beta_{\mathrm{audit}}) \le 0$, giving the stated inequality.

\emph{Part~3 (R3).} By the dual representation of total variation distance, $|\E_{P_{\mathrm{ord}}}[\tau_\pi] - \E_{P_{\mathrm{trap}}}[\tau_\pi]| \le \dTV(P_{\mathrm{ord}}, P_{\mathrm{trap}}) \le \delta_{\mathrm{ind}}$ for all measurable $\tau_\pi : \mathcal{X} \to [0, 1]$, hence $\tau_{\mathrm{trap}}(\pi) \ge \tau_{\mathrm{ord}}(\pi) - \delta_{\mathrm{ind}}$. On $P_{\mathrm{ord}}$, expected gain is at most $\tau_{\mathrm{ord}}(\pi)\Delta R$. On $P_{\mathrm{trap}}$, whenever $C > \frac{\varepsilon_{\mathrm{trap}}}{1-\varepsilon_{\mathrm{trap}}}\Delta R$, each cheating attempt yields expected payoff $\varepsilon_{\mathrm{trap}}\Delta R - (1-\varepsilon_{\mathrm{trap}})C < 0$, so expected gain on $P_{\mathrm{trap}}$ is at most $(\tau_{\mathrm{ord}}(\pi) - \delta_{\mathrm{ind}})[\varepsilon_{\mathrm{trap}}\Delta R - (1-\varepsilon_{\mathrm{trap}})C]$. Combining with mixture weights $(1-\rho)$ and $\rho$, the total expected reward $\E_{P_{\mathrm{mix}}}[\bar{R}_C(\pi)]$ is upper-bounded by
\[
  \tau_{\mathrm{ord}}(\pi)\left[(1-\rho)\Delta R + \rho\big(\varepsilon_{\mathrm{trap}}\Delta R - (1-\varepsilon_{\mathrm{trap}})C\big)\left(1 - \frac{\delta_{\mathrm{ind}}}{\tau_{\mathrm{ord}}(\pi)}\right)\right],
\]
which is strictly decreasing in $\tau_{\mathrm{ord}}(\pi)$ and strictly negative for all $\tau_{\mathrm{ord}}(\pi) \ge \tau_{\min}$ whenever $C > C_{\mathrm{R3}}^*$.
\end{proof}

\begin{remark}[Information-Theoretic Duality and Rare-Attack Floor of Regime~R3]
Total variation bounds are classically used to establish impossibility limits for \emph{monitors} trying to distinguish covert attacks from honest behavior (Neyman--Pearson power ceiling $\alpha + \dTV$). Theorem~\ref{thm:escape}(3) turns the exact same total variation inequality in favor of the \emph{defender}: if planted honeypots are close to ordinary episodes in total variation, no policy can selectively cheat on ordinary episodes while avoiding the traps. However, this guarantee requires $\tau_{\min} > \delta_{\mathrm{ind}}$ and does not block a policy that cheats at a very rare rate $\tau_{\mathrm{ord}} \le \delta_{\mathrm{ind}}$ concentrated exclusively on contexts where $P_{\mathrm{ord}} > P_{\mathrm{trap}}$.
\end{remark}

\begin{remark}[Counterexamples that delimit the claims]\label{rem:delimit}
All values are exact or recomputed by the companion script. (i)~\emph{Off-loop placement does not give non-positive covariance:} with $\beta_{\mathrm{loop}} \in \{1/10, 9/10\}$ and $\beta_{\mathrm{audit}} \in \{2/10, 8/10\}$ on two equiprobable atoms, Operator~A raises the audit survival from $1/2$ to $37/50$. (ii)~\emph{Nor is it necessary:} an audit used in the loop jointly with another detector, $\beta_{\mathrm{loop}} = \beta_{\mathrm{other}}\beta_{\mathrm{audit}}$ with $\beta_{\mathrm{audit}} = (1/2, 1)$ and $\beta_{\mathrm{other}} = (1, 1/10)$, sees its survival fall from $3/4$ to $7/12$. (iii)~\emph{Zero covariance is not sufficient under Operator~B:} $\beta_{\mathrm{loop}} \in \{1/10, 1/2, 9/10\}$, $\beta_{\mathrm{audit}} \in \{9/10, 1/10, 9/10\}$ on three equiprobable atoms have zero covariance, yet with $\eta(\Delta R + C) = 4$ the audit survival rises from $19/30$ to $0.770$, because the weight $e^{ku}$ is convex. (iv)~\emph{Collapse to cheating, not to stealth:} $p = (0.49, 0.5, 0.01)$, $\beta = (1, 1/2, 3/5)$ give a stealthy share of cheaters $3/128$ under Operator~A for every $\eta$. (v)~\emph{Enrichment can fail under Operator~B with heterogeneous rewards:} two cheaters of equal prior mass, survival $0.5$ and $0.9$, rewards $0.01$ and $10$, $C = 0$, $\eta = 0.3$: the mean survival of cheaters falls from $0.8810$ (no monitor) to $0.8748$.
\end{remark}

\section{Numerical Phase Transitions}

{\sloppy All numerical examples below are recomputed by the standalone script \texttt{verifier\_exemples\_papier\_v3.py} from the closed-form weights. The script also checks the variance identity on one example in exact rational arithmetic and runs randomized consistency checks between the stated thresholds and the directly computed weights; these catch typographical errors in formulas but are not evidence beyond the proofs.\par}

\subsection{Phase Transition across Reward Advantage and Optimization Pressure}

For representative parameters $p_{\mathrm{h}} = 0.85$, $p_{\mathrm{v}} = 0.14$, $p_{\mathrm{f}} = 0.01$, $\beta_{\mathrm{h}} = 1.0$, $\beta_{\mathrm{v}} = 0.05$, $\beta_{\mathrm{f}} = 0.75$, $\eta = 3.0$, and $C = 2.5$, Theorem~\ref{thm:selection} predicts $\Delta R_{\mathrm{A}}^* = 0.257992$, $\Delta R_{\mathrm{dom, A}}^* = 1.576778$, and $\Delta R_{\mathrm{dom, B}}^* = 2.807845$ (with asymptotic floor $0.833333$). Tables~\ref{tab:deltaR} and~\ref{tab:eta} confirm these thresholds to full precision.

\begin{table}[ht]
\centering
\caption{Phase transition across reward advantage $\Delta R$ ($p_{\mathrm{h}}=0.85, p_{\mathrm{v}}=0.14, p_{\mathrm{f}}=0.01, \beta_{\mathrm{h}}=1.0, \beta_{\mathrm{v}}=0.05, \beta_{\mathrm{f}}=0.75, \eta=3.0, C=2.5$). Bold entries highlight the crossovers $\Delta R_{\mathrm{A}}^*$, $\Delta R_{\mathrm{dom, A}}^*$, and the high-$\Delta R$ flip of Operator~B. The bold crossover rows ($0.2580$ and $1.5768$) are evaluated at the 4-decimal rounded values of $\Delta R_{\mathrm{A}}^* = 0.257992$ and $\Delta R_{\mathrm{dom, A}}^* = 1.576778$, which accounts for the $2.3 \times 10^{-5}$ difference between $\pi_{\mathrm{A}}(\Pi_{\mathrm{f}}) = 0.340917$ and $\pi_{\mathrm{A}}(\Pi_{\mathrm{h}}) = 0.340894$.}
\label{tab:deltaR}
\small
\setlength{\tabcolsep}{4pt}
\begin{tabular}{@{}rccccccc@{}}
\toprule
$\Delta R$ & $\pi_{\mathrm{no}}(\Pi_{\mathrm{f}})$ & Cond.\ no mon. & $\pi_{\mathrm{A}}(\Pi_{\mathrm{f}})$ & Cond.\ A & $\pi_{\mathrm{A}}(\Pi_{\mathrm{h}})$ & $\pi_{\mathrm{B}}(\Pi_{\mathrm{f}})$ & $\pi_{\mathrm{B}}(\Pi_{\mathrm{h}})$ \\
\midrule
$0.2000$ & $0.016221$ & $0.066667$ & $0.015593$ & $0.517241$ & $0.969854$ & $0.002821$ & $0.997043$ \\
$\mathbf{0.2580}$ & $\mathbf{0.018451}$ & $0.066667$ & $\mathbf{0.018451}$ & $0.517241$ & $0.964329$ & $0.003213$ & $0.996650$ \\
$0.8000$ & $0.044031$ & $0.066667$ & $0.081869$ & $0.517241$ & $0.841721$ & $0.010795$ & $0.989057$ \\
$1.2000$ & $0.057728$ & $0.066667$ & $0.198806$ & $0.517241$ & $0.615641$ & $0.026140$ & $0.973706$ \\
$\mathbf{1.5768}$ & $0.063492$ & $0.066667$ & $\mathbf{0.340917}$ & $0.517241$ & $\mathbf{0.340894}$ & $0.058966$ & $0.940876$ \\
$2.2000$ & $0.066157$ & $0.066667$ & $0.479040$ & $0.517241$ & $0.073856$ & $0.202970$ & $0.796883$ \\
$3.0000$ & $0.066620$ & $0.066667$ & $0.513526$ & $0.517241$ & $0.007182$ & $\mathbf{0.606385}$ & $0.393533$ \\
\bottomrule
\end{tabular}
\end{table}

\begin{table}[ht]
\centering
\caption{Asymptotic divergence as optimization pressure $\eta$ increases at fixed $\Delta R = 0.60 < \frac{1-\beta_{\mathrm{f}}}{\beta_{\mathrm{f}}}C = 0.833333$ (same base parameters as Table~\ref{tab:deltaR}).}
\label{tab:eta}
\small
\begin{tabular}{@{}rcccc@{}}
\toprule
$\eta$ & $\pi_{\mathrm{A}}(\Pi_{\mathrm{h}})$ & $\pi_{\mathrm{A}}(\Pi_{\mathrm{f}})$ & $\pi_{\mathrm{B}}(\Pi_{\mathrm{h}})$ ($C=2.5$) & $\pi_{\mathrm{B}}(\Pi_{\mathrm{f}})$ ($C=2.5$) \\
\midrule
$0.5$ & $0.977491$ & $0.011642$ & $0.941823$ & $0.010152$ \\
$1.0$ & $0.969854$ & $0.015593$ & $0.974980$ & $0.009629$ \\
$2.0$ & $0.946399$ & $0.027725$ & $0.990292$ & $0.008210$ \\
$4.0$ & $0.841721$ & $0.081869$ & $0.994178$ & $0.005808$ \\
$8.0$ & $0.325433$ & $0.348914$ & $0.997107$ & $0.002893$ \\
$16.0$ & $\mathbf{0.003955}$ & $\mathbf{0.515196}$ & $\mathbf{0.999285}$ & $\mathbf{0.000715}$ \\
\bottomrule
\end{tabular}
\end{table}

\subsection{Perfect-Stealth Niche Inversion and Escape-Regime Examples}

When a perfectly stealthy niche exists ($\beta_{\mathrm{f}} = \beta_{\mathrm{h}} = 1.0$, $\beta_{\mathrm{v}} = 0.10$, $p_{\mathrm{h}}=0.85, p_{\mathrm{v}}=0.14, p_{\mathrm{f}}=0.01$, $\eta = 2.0$, $\Delta R = 1.0$),~\eqref{eq:inversion} predicts the exact crossover $\eta(\Delta R + C) = \log(10)/0.9 = 2.558428$, i.e., $C = 0.279214$. Table~\ref{tab:niche} shows that for $C > 0.279214$, Operator~B drives absolute stealth strictly above Operator~A.

\begin{table}[ht]
\centering
\caption{Perfect-stealth niche inversion ($\beta_{\mathrm{f}} = \beta_{\mathrm{h}} = 1.0, \beta_{\mathrm{v}} = 0.10, \eta = 2.0, \Delta R = 1.0$). Above $C = 0.279214$, negative penalty amplifies absolute stealth strictly more than pure rejection. The bold crossover row ($C = 0.2792$) is evaluated at the 4-decimal rounded threshold, giving $\eta(\Delta R + C) = 2.558400$ versus the exact value $2.558428$.}
\label{tab:niche}
\small
\begin{tabular}{@{}rccccc@{}}
\toprule
Penalty $C$ & $\eta(\Delta R + C)$ & $\pi_{\mathrm{no}}(\Pi_{\mathrm{f}})$ & $\pi_{\mathrm{A}}(\Pi_{\mathrm{f}})$ & $\pi_{\mathrm{B}}(\Pi_{\mathrm{f}})$ & $\pi_{\mathrm{B}}(\Pi_{\mathrm{h}})$ \\
\midrule
$0.0000$ & $2.000000$ & $0.037731$ & $0.071924$ & $0.067487$ & $0.776336$ \\
$\mathbf{0.2792}$ & $\mathbf{2.558400}$ & $0.037731$ & $\mathbf{0.071924}$ & $\mathbf{0.071924}$ & $0.827379$ \\
$1.0000$ & $4.000000$ & $0.037731$ & $0.071924$ & $0.077603$ & $0.892711$ \\
$2.5000$ & $7.000000$ & $0.037731$ & $0.071924$ & $\mathbf{0.079814}$ & $0.918135$ \\
\bottomrule
\end{tabular}
\end{table}

For Regime~R1, with $p_{\mathrm{h}} = 0.90, p_{\mathrm{v}} = 0.09, p_{\mathrm{f}} = 0.01, \beta_{\mathrm{h}} = 1.0, \beta_{\mathrm{v}} = 0.10, \beta_{\mathrm{f}} = 0.90$, a maintenance cost $\delta_{\mathrm{maint}} = 1.25$ against $\Delta R_{\mathrm{v}} = 1.0$ yields $\Delta R_{\mathrm{net}} = -0.25$, driving honesty to $\pi_{\mathrm{B}}(\Pi_{\mathrm{h}}) = 0.999998$ at $\eta = 20, C = 2.0$. For Regime~R3, with trap fraction $\rho = 0.20$, trap false-negative rate $\varepsilon_{\mathrm{trap}} = 0.02$, indistinguishability $\delta_{\mathrm{ind}} = 0.05$, minimum targeted cheating rate $\tau_{\min} = 0.20$, and $\Delta R = 1.0$,~\eqref{eq:R3} gives $C_{\mathrm{R3}}^* = 5.462585$.

\section{Limitations and Open Problems}

Our results are intentionally derived in a static, non-parametric Gibbs equilibrium model with fixed strata: $p_{\mathrm{f}}$ and $\beta_{\mathrm{f}}$ are inputs, so the model is silent on generalisation---whether training against a detector reshapes them---which is the heart of the question that motivated it. False alarms on honest behaviour exceeding stealthy survival ($\beta_{\mathrm{f}} > \beta_{\mathrm{h}}$) are excluded by assumption. Three further open problems remain for bridging this exact characterization to practical frontier post-training:
\begin{enumerate}
  \item \textbf{Finite-time stochastic gradient dynamics on neural manifolds:} In parametric policy optimization ($\theta_{t+1} = \theta_t + \gamma_t \hat{g}_t$), implicit gradient bias and finite exploration simultaneously shape effective reachability and multi-objective maintenance costs $\delta_{\mathrm{maint}}(\theta_t)$ well before reaching any Gibbs equilibrium.
  \item \textbf{Multi-monitor coupling and cross-channel covariance:} Regime~R2 requires $\Cov_{\nu_0}(\beta_{\mathrm{loop}}, \beta_{\mathrm{audit}}) \le 0$ (necessary under both operators, sufficient under Operator~A only). When an in-loop monitor and an off-loop audit channel share latent features, characterizing worst-case degradation under constrained coupling remains open.
  \item \textbf{Dynamic monitor-versus-policy allocation and rare-attack floor:} As highlighted in AI control research~\cite{ukaisi2025}, rare interceptions create a non-stationary partial-monitoring trade-off~\cite{bartok2014} between updating the monitor and updating the policy, since policy updates can reduce future observability; moreover, Regime~R3 leaves an irreducible floor for rare attacks with $\tau_{\mathrm{ord}} \le \delta_{\mathrm{ind}}$.
\end{enumerate}
As future work on formal verification, the purely rational algebraic identity underlying Theorem~\ref{thm:selection}(2) (upon setting $E := e^{\eta \Delta R} > 0$) lends itself directly to machine-checked formalization in a proof assistant such as Lean~4.

\begin{thebibliography}{99}

\bibitem{baker2025}
B.~Baker et al.,
\emph{Monitoring Reasoning Models for Misbehavior and the Risks of Promoting Obfuscation},
arXiv:2503.11926, 2025.

\bibitem{bartok2014}
G.~Bart\'ok, D.~P.~Foster, D.~P\'al, A.~Rakhlin, and C.~Szepesv\'ari,
\emph{Partial Monitoring---Classification, Regret Bounds, and Algorithms},
arXiv:1102.2041, 2014.

\bibitem{gao2022}
L.~Gao, J.~Schulman, and J.~Hilton,
\emph{Scaling Laws for Reward Model Overoptimization},
arXiv:2210.10760, 2022.

\bibitem{greenblatt2023}
R.~Greenblatt et al.,
\emph{AI Control: Improving Safety Despite Intentional Subversion},
arXiv:2312.06942, 2023.

\bibitem{greenblatt2024}
R.~Greenblatt et al.,
\emph{Alignment Faking in Large Language Models},
arXiv:2412.14093, 2024.

\bibitem{griffin2024}
C.~Griffin et al.,
\emph{Games for AI Control: Models of Safety Evaluations of AI Deployment Protocols},
arXiv:2409.07985, 2024.

\bibitem{halberstadt2026}
Y.~Halberstadt,
\emph{Post training is fundamentally dangerous},
LessWrong, post \texttt{CoRJ9saddvsXti4m2}, November 22, 2026.

\bibitem{hardt2016}
M.~Hardt, N.~Megiddo, C.~Papadimitriou, and M.~Wootters,
\emph{Strategic Classification},
arXiv:1506.06980, 2016.

\bibitem{hubinger2024}
E.~Hubinger et al.,
\emph{Sleeper Agents: Training Deceptive LLMs that Persist Through Safety Training},
arXiv:2401.05566, 2024.

\bibitem{joagdev1983}
K.~Joag-Dev and F.~Proschan,
\emph{Negative Association of Random Variables with Applications},
The Annals of Statistics, 11(1):286--295, 1983.

\bibitem{skalse2022}
J.~Skalse, N.~H.~R.~Howe, D.~Krasheninnikov, and D.~Krueger,
\emph{Defining and Characterizing Reward Gaming},
arXiv:2209.13085, 2022.

\bibitem{terekhov2025}
M.~Terekhov et al.,
\emph{Control Tax: The Price of Keeping AI in Check},
arXiv:2506.05296, 2025.

\bibitem{testingthewaters2026}
\texttt{testingthewaters},
Comment on \emph{Post training is fundamentally dangerous},
LessWrong, comment \texttt{gFufuPsjYMbutMtkb}, November 2026.

\bibitem{ukaisi2025}
UK AI Security Institute (UK AISI),
\emph{Research Areas in AI Control},
LessWrong, post \texttt{rGcg4XDPDzBFuqNJz}, August 2025.

\bibitem{zhuang2021}
S.~Zhuang and D.~Hadfield-Menell,
\emph{Consequences of Misaligned AI},
arXiv:2102.03896, 2021.

\bibitem{fisher1930}
R.~A.~Fisher,
\emph{The Genetical Theory of Natural Selection},
Clarendon Press, Oxford, 1930.

\bibitem{korbak2022}
T.~Korbak, E.~Perez, and C.~L.~Buckley,
\emph{RL with KL Penalties is Better Viewed as Bayesian Inference},
Findings of EMNLP 2022, arXiv:2205.11275.

\bibitem{korbak2025}
T.~Korbak et al.,
\emph{Chain of Thought Monitorability: A New and Fragile Opportunity for AI Safety},
arXiv:2507.11473, 2025.

\bibitem{price1970}
G.~R.~Price,
\emph{Selection and Covariance},
Nature, 227:520--521, 1970.

\bibitem{robertson1966}
A.~Robertson,
\emph{A Mathematical Model of the Culling Process in Dairy Cattle},
Animal Production, 8:95--108, 1966.

\end{thebibliography}

\end{document}
